\documentclass[12pt,a4paper]{article}
\usepackage[utf8]{inputenc}
\usepackage[T1]{fontenc}
\usepackage{geometry}
\geometry{margin=2.5cm}
\usepackage{setspace}
\onehalfspacing
\usepackage{hyperref}

\begin{document}

\section{Literature Review}

\bigskip
\noindent\textbf{The Anti-Money Laundering Problem}

\noindent Money laundering remains one of the most pervasive challenges facing the global financial system. The United Nations Office on Drugs and Crime estimates that between two and five percent of global GDP is laundered annually, yet only a fraction of illicit flows are ever detected (Cheng, Zou, Xiang, \& Jiang, 2024). Banks are legally obligated to monitor transactions and file suspicious activity reports (SARs), but the rule-based systems that underpin most transaction monitoring frameworks are notoriously imprecise. Gerlings and Constantiou (2023), in a two-year action research study conducted with a European bank, found that traditional rule-based scenarios generate false positive rates between 95 and 98 percent, meaning that the vast majority of alerts reviewed by investigators are ultimately unproductive. Cardoso, Saleiro, and Bizarro (2022) corroborate this figure, noting that legacy AML systems routinely exceed 95 percent false positive rates.

The operational cost of these false positives is substantial: investigators must manually review each flagged alert within tight regulatory timeframes, and the sheer volume of unproductive alerts diverts resources away from genuinely suspicious activity. Vallarino (2025) estimates that financial fraud causes global losses exceeding 1.8 trillion dollars annually, with traditional systems unable to keep pace with increasingly sophisticated laundering schemes. This inefficiency has motivated a sustained effort to apply machine learning methods to AML, first through tabular approaches and more recently through graph-based architectures that can capture the relational structure inherent in financial transaction data.

\bigskip
\noindent\textbf{Traditional Machine Learning for AML}

\noindent Early applications of machine learning to AML and fraud detection relied on supervised tabular models. Random forests, gradient-boosted trees such as XGBoost, and logistic regression have been widely deployed for transaction-level classification, typically operating on hand-engineered features derived from transaction amounts, frequencies, counterparty statistics, and temporal patterns (Cheng et al., 2024). Bathala, Sasikala, Prasad, Begum, Mahajan, and Sreedevi (2024) demonstrated the use of deep feedforward neural networks with SHAP-based explainability on synthetic transaction data, achieving weighted F1 scores around 0.91 but acknowledging that synthetic datasets may not capture real-world complexity.

A fundamental limitation of all tabular approaches, however, is that they treat each transaction or customer as an independent observation, discarding the rich relational context in which financial activity occurs. A single suspicious transaction may appear benign in isolation but reveal its illicit nature only when examined alongside the network of accounts, counterparties, and intermediary entities through which funds have flowed. This structural blindness---the inability to reason about the topology of financial relationships---has driven the field toward graph-based representations.

\bigskip
\noindent\textbf{From Tabular Data to Graph Representations}

\noindent The recognition that financial crime is fundamentally a relational phenomenon has motivated a shift toward graph-based representations of transaction data. As Johannessen and Jullum (2023) argue, money laundering is a social activity in which groups of organized criminals collaborate to conceal illicit proceeds---and the patterns that betray this collaboration are often invisible when customers or transactions are examined in isolation. Network analysis can uncover circumstances that would be impossible to detect through a purely entity-based approach, simply because the necessary relational information is not present in tabular feature sets. At the simplest level, accounts can be modeled as nodes and transactions as directed edges, forming a homogeneous transaction graph. However, this representation discards critical distinctions between different types of entities and relationships.

Real banking environments are inherently heterogeneous. A bank's data contains multiple entity types---individual customers, corporate customers, external accounts, merchants, devices, geographic locations---connected through diverse relationship types such as financial transfers, shared ownership, business ties, and common attributes. Singh, Prasad, Michael, Kaphungkui, and Singh (2024) emphasize that financial data, especially involving credit cards, encompasses diverse entities like cardholders, merchants, and transactions, each with distinct attributes and relationships that homogeneous graphs cannot fully capture. A heterogeneous graph, by contrast, preserves these distinctions: nodes carry type-specific feature sets, and edges encode both the nature and the properties of each relationship. This matters concretely for AML because a laundering scheme might involve a chain of shell companies connected through ownership edges, with layering transactions flowing between their accounts, and the suspicious pattern only becomes visible when these different relation types are jointly considered. Johannessen and Jullum (2023) further note that existing heterogeneous GNN methods such as RGCN (Schlichtkrull et al., 2018), HAN (Wang et al., 2019), and HGT (Hu et al., 2020) were not designed to incorporate edge features---yet transaction properties like amounts, timestamps, and currency types carry essential information for an effective AML learning task.

Yu and colleagues (2024) explored knowledge graph construction for AML, representing bank accounts as nodes and transactions as edges with attributes including amount, currency, and payment format. They applied Graph Attention Networks (GAT) with focal loss to address class imbalance on a synthetic IBM dataset of 1,243 transactions, finding that timestamp and payment format were the most discriminative features. However, performance degraded as test set size increased---accuracy dropped from 0.70 to 0.65 when scaling from 2,000 to 4,000 records---an important cautionary note about generalization on small-scale synthetic data. The question of how to construct and learn from financial graphs at realistic scale, while preserving the heterogeneous structure and edge features that make graph representations valuable in the first place, has thus become a central concern in the literature.

Once a graph representation is constructed, there are broadly two strategies for incorporating its information into a detection model. The first is to \emph{explicitly} compute graph-derived features---such as centrality metrics, random walk statistics, or neighborhood aggregates---and feed them as additional inputs to a conventional classifier. Naser Eddin et al. (2021) exemplify this approach, extracting degree statistics and GuiltyWalker propagation scores from a dynamic transaction graph and using them alongside entity-level features in a LightGBM model. While effective, such hand-crafted features may not capture the structural patterns most informative for the downstream task, and their design requires considerable domain expertise. The second strategy is to \emph{implicitly} learn representations through graph neural networks that operate directly on the graph structure, jointly optimizing feature extraction and classification in an end-to-end fashion. As Cardoso et al. (2022) note, implicitly derived representations through node embedding approaches have generally achieved better downstream classification results than curated graph features, supporting the argument that end-to-end learned representations contain additional, meaningful information beyond what manual feature engineering can capture.

\bigskip
\noindent\textbf{Graph Neural Networks for Fraud Detection}

\noindent The application of graph neural networks to financial fraud detection has grown rapidly in recent years. Cheng et al. (2024) provide a comprehensive review of over 100 studies, organizing GNN approaches for fraud detection into four categories: graph convolutional networks (GCNs), graph attention networks (GATs), graph temporal networks, and heterogeneous graph neural networks. GCNs, introduced by Kipf and Welling (2017), aggregate neighborhood information through spectral convolutions on the graph Laplacian, while GraphSAGE (Hamilton, Ying, \& Leskovec, 2017) extends this to an inductive setting through neighborhood sampling and aggregation. These homogeneous architectures laid the groundwork for graph-based fraud detection but treat all nodes and edges as identical, which is a poor fit for the multi-relational structure of financial networks. Cheng et al. further note that GCNs suffer from over-smoothing as depth increases, losing the distinct node characteristics that are essential for distinguishing illicit from legitimate behavior.

\bigskip
\noindent\textbf{Heterogeneous Graph Neural Network Architectures}

\noindent Recognizing the limitations of homogeneous approaches, a line of work has developed GNN architectures that explicitly model heterogeneous node and edge types. The Heterogeneous Graph Transformer (HGT), proposed by Hu, Dong, Wang, and Sun (2020), introduced node-type and edge-type dependent attention mechanisms for heterogeneous graphs. HGT uses separate key, query, and value projections parameterized by meta-relation triplets---source node type, edge type, and target node type---allowing the model to learn type-specific importance weights during message passing. On the Open Academic Graph (179 million nodes, 2 billion edges), HGT outperformed prior methods by 9 to 21 percent across downstream tasks, establishing the viability of heterogeneous attention at web scale.

Rao et al. (2022) built directly on HGT in developing xFraud, an explainable fraud detection system deployed at eBay. The xFraud detector constructs a heterogeneous graph with transaction nodes linked to buyers, sellers, payment tokens, email addresses, and shipping addresses, then applies heterogeneous mutual attention and message passing with self-attention mechanisms. On eBay's industrial-scale dataset (1.1 billion nodes, 3.7 billion edges), xFraud achieved an AUC of 0.91, outperforming both vanilla GCN-based approaches (GEM, AUC 0.90) and standard GAT (AUC 0.89). Notably, xFraud replaced HGT's original HGSampling with a GraphSAGE-style sampler for efficiency on sparse financial graphs (2.12 edges per node, versus 11.17 in the academic graph), reducing inference time while maintaining detection quality. The system also introduced a hybrid explainer combining task-aware GNNExplainer edge weights with task-agnostic centrality measures, demonstrating that neither approach alone dominates across all test communities---a finding with direct implications for explainability requirements in AML, discussed further below.

\bigskip
\noindent\textbf{Heterogeneous GNNs for Anti-Money Laundering}

\noindent Within the specific domain of anti-money laundering, several studies have explored heterogeneous GNN architectures on transaction data at varying scales and with different modeling choices. Johannessen and Jullum (2023) presented the first published application of GNNs to large-scale real-world heterogeneous bank data from DNB, Norway's largest financial institution. Their Heterogeneous Message Passing Neural Network (HMPNN) extends standard MPNNs by introducing edge-type-specific message functions that incorporate edge features---a critical design choice, since transaction properties such as amounts, timestamps, and currency types carry essential information for AML that purely structural approaches would miss. On a graph of over 5 million nodes and 10 million edges spanning individual customers, organization customers, and external accounts connected through financial transactions and business ties, HMPNN demonstrated that heterogeneous graph structure preserves information that homogeneous approaches lose. This work is particularly significant because it validates heterogeneous GNNs on genuine banking data rather than synthetic benchmarks, though the proprietary nature of the data limits reproducibility.

Naser Eddin, Bono, Aparício, Polido, Ascensão, and Bizarro (2021) approached the problem from the perspective of alert optimization rather than raw transaction classification. They combined entity-centric profile features (computed across five time windows ranging from one day to two months) with graph-based features including node degree statistics and GuiltyWalker propagation scores, constructed through sliding-window dynamic graphs. Their LightGBM-based system achieved an 80 percent reduction in false positives while maintaining over 90 percent true positive detection on a real banking dataset of approximately 500,000 alerted transfers. Graph-based features alone contributed a 13.4 percentage point improvement in recall, underscoring the practical value of graph-augmented approaches for operational AML.

Several other contributions have extended the heterogeneous GNN paradigm for financial crime in complementary directions. Singh, Prasad, Michael, Kaphungkui, and Singh (2024) proposed a heterogeneous graph autoencoder for credit card fraud detection, constructing a tripartite graph of cardholders, merchants, and transactions. Their autoencoder, trained exclusively on legitimate transactions, flags anomalous reconstruction errors as potential fraud, achieving an AUC-PR of 0.89 and an F1 score of 0.81. The autoencoder paradigm is notable because it sidesteps the label scarcity problem by learning what constitutes normal behavior rather than requiring explicit fraud labels. Stan (2024), in an MSc thesis at Utrecht University, applied provably powerful directed multigraph GNNs to AML on a large-scale synthetic transaction dataset of approximately 9.5 million transactions with a 0.14 percent illicit rate. Stan introduced ego IDs, heterogeneous message passing, and port numbering to enhance expressiveness beyond what standard message-passing GNNs can capture. Shwetha, Shreyas, Siddarameshwara, Srinidhi, and Srujan (2025) took a different angle with group-aware deep graph learning, using a community-centric encoder that aggregates flagged transactions into ``super-nodes'' to detect organized gang-level laundering rather than individual suspicious transactions, tested on data from one of the world's largest bank card alliances.

\bigskip
\noindent\textbf{Self-Supervised Graph Learning for AML}

\noindent While the studies reviewed above demonstrate the value of heterogeneous graph structure for AML, they share a common dependence on labeled data---a resource that is particularly scarce and unreliable in the anti-money laundering domain. Ground truth labels in AML are inherently incomplete: they reflect only detected and confirmed laundering, not the true underlying illicit activity. This means that supervised models are trained on a biased and incomplete view of the positive class, and the negative class (ostensibly legitimate transactions) is contaminated by undetected laundering. A parallel line of research has therefore explored self-supervised graph learning as a strategy for overcoming this severe label scarcity.

Early work on graph-based AML already hinted at the value of learned representations. As noted by Cardoso et al. (2022), Weber et al. (2019) compared GCN-based models against a random forest on Bitcoin transaction data and found that the GCN actually performed worse when input features were already well-curated---yet combining GCN-derived embeddings with those features improved upon all baselines. This finding suggests that graph representations capture complementary structural information even when they do not dominate in isolation. Lo et al. (2022) extended this line of work by applying the Deep Graph Infomax (DGI) self-supervised objective to derive additional input features for a supervised classifier, further improving results. In their work, however, self-supervision served as a stepping stone to enhance supervised performance rather than as the primary detection mechanism.

Cardoso et al. (2022) introduced LaundroGraph, the first fully self-supervised AML system. Operating on a customer-transaction bipartite graph from a real European bank (3.2 million customer nodes and 17.7 million transaction nodes), LaundroGraph uses a three-layer GAT encoder with self-supervised link prediction to learn node embeddings, then derives anomaly scores from reconstruction errors on held-out transactions. LaundroGraph outperformed non-graph baselines by 12 percentage points in AUC (94.83 percent versus 82.58 percent for LightGBM), and also surpassed the contrastive self-supervised baseline Deep Graph Infomax by nearly 9 percentage points. This result established that self-supervised graph learning can meaningfully improve AML detection even in the absence of labeled data. However, LaundroGraph operates on a homogeneous bipartite graph with only two node types and a single edge type, leaving open the question of whether richer heterogeneous graph structures could further improve self-supervised representations.

The Heterogeneous Graph Masked Autoencoder (HGMAE), proposed by Tian, Dong, Zhang, Zhang, and Chawla (2022) and presented as an oral paper at AAAI 2023, represents the state of the art in self-supervised pretraining on heterogeneous graphs. HGMAE extends the masked autoencoding paradigm to heterogeneous settings through three complementary training objectives: metapath-based edge reconstruction, which reconstructs masked edges in metapath-derived adjacency matrices with semantic-level attention weighting; target attribute restoration, which recovers masked node features using learnable mask tokens with an adaptive masking schedule that increases difficulty during training; and positional feature prediction, which captures structural roles through metapath2vec-based positional encodings. The encoder is based on the Heterogeneous Attention Network (HAN), employing hierarchical attention---first at the node level within each meta-path, then at the semantic level across meta-paths---to learn type-aware representations. On standard heterogeneous graph benchmarks, HGMAE achieves 92.71 percent Micro-F1 on DBLP and 90.24 percent on ACM for node classification, competitive with or superior to fully supervised methods. However, HGMAE has not been applied to financial transaction data or AML detection, leaving open the question of whether its self-supervised heterogeneous representations transfer effectively to this domain.

Paladugu (2025) explored a related but distinct direction, combining GNN-generated graph embeddings with homomorphic encryption for privacy-preserving AML. This work addresses the cross-border collaboration challenge---where banks in different jurisdictions cannot share raw transaction data due to privacy regulations---but focuses primarily on the privacy dimension rather than detection performance.

\bigskip
\noindent\textbf{Temporal Dynamics and Continual Learning}

\noindent Beyond the question of how to learn effective representations from transaction graphs, a further challenge concerns the temporal dynamics of money laundering itself. Laundering schemes evolve over time as criminals adapt to detection mechanisms, creating concept drift that degrades model performance. Deprez, Wei, Verbeke, Baesens, Mets, and Verdonck (2025) provided a comprehensive review of continual graph learning methods for AML, evaluating replay, regularization, and architecture-based strategies to handle concept drift without catastrophic forgetting. Their work highlights a critical operational concern: AML models must not only detect current laundering patterns but also retain knowledge of previously observed tactics, since criminals may revert to older methods once they believe a particular scheme is no longer under surveillance. Among the methods reviewed, architecture-based approaches such as PackNet showed the strongest overall performance, though replay-based methods remained competitive at the cost of storing historical data---a potential privacy concern in the banking context.

The hybrid LSTM-GraphSAGE architecture explored in recent work (International American University, 2025) represents an alternative approach to temporal modeling, combining LSTM-based temporal sequence modeling with GraphSAGE-based relational modeling in a joint architecture. On a large-scale synthetic AML dataset, this hybrid model achieved 95.4 percent accuracy and an AUC of 0.97, outperforming both its constituent components in isolation (LSTM alone: 91.5 percent accuracy; GraphSAGE alone: 92.8 percent) as well as traditional baselines including Random Forest, Na\"ive Bayes, and XGBoost. While promising, the practical applicability of such hybrid architectures to real banking environments at production scale remains to be validated.

\bigskip
\noindent\textbf{Evaluation Metrics Under Extreme Class Imbalance}

\noindent The choice of evaluation metrics in AML research deserves careful consideration, given the extreme class imbalance typical of transaction monitoring data. Real-world and synthetic AML datasets commonly exhibit illicit rates well below one percent---often as low as 0.1 to 0.2 percent---meaning that a naive classifier labeling all transactions as legitimate can achieve over 99 percent accuracy. Accuracy is therefore a misleading metric, and the field has converged on metrics that are sensitive to performance on the minority class.

The area under the receiver operating characteristic curve (AUC-ROC) measures discriminative ability across all classification thresholds and is widely reported: xFraud achieved 0.91 on eBay data (Rao et al., 2022), while LaundroGraph reached 0.95 on real banking data (Cardoso et al., 2022). However, AUC-ROC can be overly optimistic under severe imbalance because it credits performance at high false positive rates that would be operationally unacceptable. The area under the precision-recall curve (AUC-PR) is arguably more informative in these settings, as it directly measures the tradeoff between the fraction of flagged transactions that are truly illicit (precision) and the fraction of all illicit transactions that are detected (recall). Singh et al. (2024) reported AUC-PR of 0.89 on credit card fraud, while Stan (2024) evaluated on large-scale synthetic AML data using similar precision-recall metrics. The F1 score, the harmonic mean of precision and recall, provides a single-number summary of this tradeoff and is reported across most studies in the field.

From an operational perspective, metrics beyond statistical performance also matter. Naser Eddin et al. (2021) emphasized false positive reduction as a key metric---their system reduced false positives by 80 percent---because every false positive represents wasted investigator time. Gerlings and Constantiou (2023) further noted that time-to-SAR (the time from alert generation to filing a suspicious activity report) is an operationally critical metric that is rarely reported in the academic literature but directly affects regulatory compliance. A model that achieves high AUC but generates so many alerts that investigators cannot process them within regulatory deadlines provides little practical value.

\bigskip
\noindent\textbf{Explainability and Regulatory Requirements}

\noindent The question of explainability in AML deserves particular emphasis given the regulatory environment in which these models operate. Banks must be able to justify their monitoring decisions to auditors and financial supervisory authorities, and a model that improves detection but cannot explain its reasoning may be unacceptable from a compliance perspective. Vallarino (2025) surveyed compliance challenges for GNN-based fraud detection, noting that regulatory frameworks such as GDPR and the Bank Secrecy Act require not just accurate detection but also transparent decision-making processes.

Gerlings and Constantiou (2023) documented a fundamental disconnect between the explanations that machine learning models provide and the explanations that AML investigators need. Their study found that when a bank deployed an augmentation model providing risk scores to investigators, the scores introduced cognitive bias without providing sufficient context for why a transaction was flagged. Investigators desired explanations grounded in the business domain---which counterparties were involved, what transaction patterns were unusual, how the flagged activity related to known typologies---rather than statistical performance summaries such as ROC curves or precision scores. Rao et al. (2022) addressed this partially through xFraud's hybrid explainer, which combines GNNExplainer edge weights with graph centrality measures to identify the most informative edges in a prediction's subgraph. Their evaluation across 41 annotated communities showed that the hybrid approach outperformed either component alone, achieving top-5 hit rates of up to 0.47.

The tension between model complexity and interpretability remains unresolved, though heterogeneous GNNs offer a structural advantage: attention weights over different relation types can reveal which kinds of relationships---shared counterparties, unusual transaction chains, atypical geographic patterns---contributed most to a prediction, providing more topology-aware explanations than per-feature importance scores from tabular models. This structural interpretability is particularly relevant for AML, where investigators need to trace specific transaction paths rather than understand aggregate feature contributions.

% TODO(human): Write 5-8 lines positioning this thesis's contribution.
% Consider: (1) No prior work applies HGMAE to AML/financial crime.
% (2) Self-supervised + heterogeneous is underexplored in this domain.
% (3) Your thesis combines HGMAE pretraining with downstream HMPNN/HGT classification.
% (4) You address both label scarcity (H2) and structural richness (H1).
% (5) Frame how this fills the gap between LaundroGraph (self-supervised but homogeneous bipartite)
%     and Johannessen & Jullum (heterogeneous but fully supervised).

\end{document}
